\label{cap:trabalhosfuturos}

\begin{itemize}
    \item Realizar estudos comparativos com medidas obtidas por DXA, permitindo a validação
    experimental dos modelos biométricos e a investigação de estratégias de calibração
    individualizada.

    \item Ampliar a coleta de dados em estudos de campo com a balança CF577BLE, aumentando o
    número de participantes e o volume de medições por usuário, o que possibilita análises
    estatísticas mais robustas e modelos mistos mais avançados.

    \item Investigar técnicas de \emph{fine-tuning} supervisionado e bases sintéticas
    especializadas para ampliar a capacidade de interpretação da assistente em cenários
    específicos de composição corporal.

    \item Expandir a compatibilidade do sistema para outras balanças e dispositivos de
    bioimpedância compatíveis com Bluetooth Low Energy (BLE).

    \item Evoluir a aplicação web progressiva (PWA) para versões móveis dedicadas,
    incorporando funcionalidades adicionais de acompanhamento da saúde e integração com
    dispositivos vestíveis.

    \item Avaliar modelos computacionais mais sofisticados para contextualização e correção
    das medições, considerando fatores como hidratação, alimentação e atividade física.

    \item Conduzir o estudo de campo com usuários voluntários. Nesse estudo, recrutar uma
    amostra de conveniência de cerca de 15 a 30 participantes maiores de 18 anos, com
    condições físicas para usar a balança e mediante assinatura do Termo de Consentimento
    Livre e Esclarecido (TCLE), excluindo quem tiver limitações que inviabilizem a medição ou
    não concordar com o uso dos dados.

    \item Executar o procedimento de coleta previsto: cadastro de cada participante (sexo,
    idade e altura), leitura dos dados brutos (peso e impedância) pela API Web Bluetooth e
    registro dos parâmetros de composição corporal, além de informações contextuais opcionais
    como hábitos de atividade física e objetivos de saúde.

    \item Aplicar a escala \emph{System Usability Scale} (SUS) com usuários e submeter as
    interpretações geradas pela assistente Vida à validação por especialistas das áreas de
    saúde, nutrição e educação física.

    \item Desenvolver e validar ajustes por regressão (linear e multivariada) específicos
    para a população brasileira e estratégias adicionais de compatibilização entre
    frequências, complementando a interpolação para 50~kHz já implementada.
\end{itemize}
